% فهرست علائم اختصاری
\chapter*{فهرست علائم اختصاری}
\phantomsection
\addcontentsline{toc}{chapter}{فهرست علائم اختصاری}

در این پایان‌نامه به‌منظور سهولت و کوتاه‌تر شدن متن، از علائم اختصاری زیر
استفاده شده است:

\begin{center}
\begin{tabular}{|c|c|}
\hline
\textbf{علامت اختصاری} & \textbf{شرح} \\
\hline
\lr{PINN} & شبکه عصبی آگاه از فیزیک \\
\hline
\lr{PDE} & معادله دیفرانسیل جزئی \\
\hline
\lr{ODE} & معادله دیفرانسیل معمولی \\
\hline
\lr{ANN} & شبکه عصبی مصنوعی \\
\hline
\lr{MLP} & پرسپترون چندلایه \\
\hline
\lr{MSE} & میانگین مربعات خطا \\
\hline
\lr{SSE} & مجموع مربعات خطا \\
\hline
\lr{IC} & شرط اولیه \\
\hline
\lr{BC} & شرط مرزی \\
\hline
\lr{RK} & رانگ ـ کوتا \\
\hline
\lr{L-BFGS} & روش شبه‌نیوتنی حافظه محدود \\
\hline
\lr{ReLU} & واحد خطی یکسوشده \\
\hline
\lr{GPU} & واحد پردازش گرافیکی \\
\hline
\lr{CPU} & واحد پردازش مرکزی \\
\hline
\end{tabular}
\end{center}
